
\section{Remote Sensing and Satellite Technology Overview}
Remote sensing represents the foundational science and technological methodology of acquiring information regarding the physical, chemical, and biological attributes of the Earth's surface without direct in-situ physical contact \cite{R201}. This objective collection of spatial information is achieved by measuring, recording, and analyzing electromagnetic radiation (EMR) that is reflected, emitted, or backscattered from terrestrial features and atmospheric components \cite{R211}. The primary operational platforms for civilian remote sensing are spaceborne satellites orbiting the Earth in various configurations, including Sun-synchronous polar orbits and geostationary orbits. These spaceborne platforms house specialized payloads categorized into two principal sensor modalities: passive and active systems \cite{R212}.

\subsection{Passive Sensors (Optical and Thermal)}
Passive remote sensing systems rely on external, natural sources of energy to illuminate the Earth's surface, utilizing solar radiation as the primary source of illumination. Sensors mounted on prominent Earth observation platforms, such as the Multi-Spectral Instrument (MSI) aboard the European Space Agency's Sentinel-2 constellation or the Operational Land Imager (OLI) aboard the NASA/USGS Landsat-8 and Landsat-9 satellites, capture and digitize this reflected solar energy across discrete, pre-defined bands of the electromagnetic spectrum \cite{R203, R213}. 

The optical spectrum utilized in civil remote sensing spans several vital spectral regions. The visible spectrum, comprised of the Blue (450–515 nm), Green (525–600 nm), and Red (630–690 nm) channels, records the primary wavelengths of light, which are critical for natural color rendering, visual interpretation, and elementary land-cover classification \cite{R214}. Beyond the visible range, the Near-Infrared (NIR) band (700–1100 nm) is highly sensitive to the internal cellular structure of healthy vegetation. Specifically, the spongy mesophyll layer of leaves exhibits a pronounced reflectance of NIR energy, creating a sharp spectral rise known as the "infrared shoulder," which serves as the physical basis for vegetative biomass estimation \cite{R215}. Further along the spectrum, the Shortwave-Infrared (SWIR) bands (1400–3000 nm) respond dynamically to leaf water content, soil moisture dynamics, and the mineralogical composition of exposed rocks, making them essential for separating bare soil from artificial, man-made structures \cite{R216}. Finally, the Thermal Infrared (TIR) region (8000–14000 nm) measures the longwave thermal energy emitted directly from the Earth's surface as a function of its temperature, facilitating urban heat island mapping, thermal plume detection from industrial installations, and surface temperature estimation \cite{R217}.

The physical interaction of solar radiation with surface materials is mathematically characterized by spectral reflectance curves, which represent the percentage of incident energy reflected by a target as a function of wavelength \cite{R218}. Healthy vegetation absorbs red light via chlorophyll pigments and reflects NIR via cellular structures, whereas bare soil exhibits a monotonically increasing reflectance curve across the visible and infrared bands. Liquid water absorbs almost all NIR and SWIR radiation, showing high contrast against dry land. However, optical systems are heavily constrained by atmospheric scattering (Rayleigh and Mie scattering) and cloud cover, which completely obscure terrestrial features in the visible and infrared bands \cite{R219}.

\subsection{Active Sensors (Radar/SAR)}
Active remote sensing systems, most notably Synthetic Aperture Radar (SAR) systems aboard the Sentinel-1 constellation, mitigate the limitations of passive sensors by providing their own illumination source. These systems emit microwave pulses toward the Earth's surface and measure the intensity and phase of the backscattered echo \cite{R220}. SAR systems operate at centimeter-scale wavelengths, such as the C-band (~5.6 cm), which can penetrate clouds, smoke, dust, and precipitation, enabling continuous day-and-night observation \cite{R221}.

The intensity of the SAR backscatter is governed by two principal physical properties: the dielectric constant of the target material, which is highly sensitive to moisture content, and the surface roughness relative to the radar wavelength. Surface scatter can be categorized into three distinct scattering mechanisms. Specular reflection occurs when smooth surfaces, such as calm water bodies, asphalt, or concrete pavements, reflect the radar pulse away from the sensor, appearing dark in SAR imagery \cite{R222}. Volume scattering arises when complex, multi-layered structures, such as dense forest canopies or crops, scatter the signal in multiple directions, returning a moderate response \cite{R223}. Double-bounce scattering occurs when orthogonal surfaces, such as buildings, shipping containers, and metallic vessels, act as corner reflectors, returning a highly intense signal back to the sensor \cite{R224}. By analyzing polarization configurations, such as VV (Vertical-Vertical) and VH (Vertical-Horizontal) backscatter, SAR pipelines can distinguish structural changes in land cover (e.g., deforestation or road construction) even in regions with persistent cloud cover \cite{R225}.

---

\section{Economic Intelligence from Space: Historical Context and Proxies}
The use of Earth Observation (EO) data to analyze human economic activity, known as "Space-Based Economic Intelligence," has grown rapidly over the last two decades. Historically, remote sensing was restricted to military and national intelligence agencies. However, the democratization of open satellite data and cloud computing has allowed academic researchers and financial institutions to extract micro- and macroeconomic signals directly from orbital data \cite{R226}.

\subsection{Night-time Lights (NTL) as a Macroeconomic Proxy}
The foundational milestone in satellite economics was established by Henderson, Storeygard, and Weil (2012), who demonstrated that changes in nocturnal light emissions (captured by DMSP-OLS and later Suomi-NPP VIIRS) correlate strongly with regional and national Gross Domestic Product (GDP) growth \cite{R227}. NTL radiance serves as a powerful proxy for economic activity because light consumption is highly correlated with electricity use, urbanization, and industrial manufacturing \cite{R228}.

Subsequent work by Chen and Nordhaus (2011) showed that in countries with weak statistical reporting agencies, NTL data can refine GDP estimates by eliminating administrative biases and reporting errors \cite{R229}. The transition from DMSP-OLS to the Visible Infrared Imaging Radiometer Suite (VIIRS) Day/Night Band (DNB) in 2013 provided higher spatial resolution (500m vs 2.7km) and eliminated sensor saturation, permitting the quantitative assessment of local industrial parks and urban corridors \cite{R230, R231}.

\subsection{Microeconomic and Supply Chain Proxies}
Beyond aggregate GDP estimation, researchers have developed spatial indicators to track specific economic sectors. Industrial tracking involves monitoring oil stockpiles by measuring the shadows cast inside floating-roof oil storage tanks \cite{R232}, and tracking shipping containers and vessel movements in deep-sea ports \cite{R233}. Agricultural monitoring uses NDVI time series to estimate crop yields, map agricultural boundaries, and forecast supply disruptions \cite{R234}. In extractive economies, remote sensing tracks mineral and fossil fuel extraction. Satellite remote sensing is particularly valuable for identifying informal or illegal mining, where operations are unmapped and hidden from official registries \cite{R235}.

---

\section{Open Satellite Data Ecosystems}
The modern geospatial data pipeline relies on open-access platforms that distribute processed satellite data. The transition from manual file transfers to programmatic REST APIs has transformed the geospatial engineering domain:

\begin{itemize}
  \item \textbf{Google Earth Engine (GEE)}: Launched in 2010, GEE revolutionized geospatial analytics by co-locating a petabyte-scale archive of public satellite data (Sentinel, Landsat, MODIS) with a planetary-scale parallel computing engine \cite{R207}. GEE uses a map-reduce architecture to execute complex spatial queries on Google's cloud clusters, bypassing local RAM and storage constraints \cite{R236}.
  \item \textbf{Copernicus Data Space Ecosystem (CDSE)}: Deployed in 2023 by the European Commission and ESA, CDSE replaced the legacy Copernicus Open Access Hub (SciHub) \cite{R202}. CDSE standardizes data distribution through OData, Sentinel Hub, and openEO APIs, providing advanced cloud-masking and processing-on-the-fly capabilities \cite{R237}.
  \item \textbf{NASA Earthdata \& USGS EarthExplorer}: Provide access to American civilian satellite archives, distributing Landsat datasets and VIIRS low-light radiance products \cite{R205, R238}.

\end{itemize}
---

\section{Anomaly Detection in Geospatial Data}
Geospatial anomaly detection involves isolating pixels or regions whose spectral or radar signatures deviate from historical baselines or spatial neighborhoods. These anomalies are classified into three types \cite{R239}:
\begin{enumerate}
  \item \textbf{Point Anomalies}: Single pixels with extreme values (e.g., thermal anomalies indicating gas flaring or volcanic activity) \cite{R240}.
  \item \textbf{Contextual Anomalies}: Values that appear normal in isolation but are anomalous within their spatial or temporal context (e.g., low NDVI in a protected rainforest area) \cite{R241}.
  \item \textbf{Collective/Temporal Anomalies}: Sequences of data points indicating a structural shift in land cover (e.g., a sharp increase in bare soil index marking a new mining pit) \cite{R242}.

\end{enumerate}
Detecting these anomalies is complicated by seasonal vegetation changes (phenology), soil moisture shifts, and sensor degradation over time \cite{R243}. Traditional thresholding methods (e.g., flagging pixels where an index drops more than two standard deviations below the mean) often suffer from high false-alarm rates, highlighting the need for multivariate statistical models \cite{R244}.

---

\section{Machine Learning for Earth Observation}
Machine learning algorithms have become central to analyzing Earth Observation datasets, falling into three primary approaches:

\begin{itemize}
  \item \textbf{Supervised Learning}: Random Forests, Support Vector Machines (SVM), and Convolutional Networks (CNNs) achieve high accuracy in land-cover classification when trained on large labeled datasets \cite{R245}. However, acquiring high-quality ground-truth labels for anomalies (such as illegal mines or unmapped industrial expansions) is difficult, which limits the utility of supervised methods \cite{R246}.
  \item \textbf{Unsupervised Learning}: Unsupervised models, such as Isolation Forests (Liu et al., 2008), isolate outliers by randomly partitioning the feature space \cite{R206}. Because anomalies require fewer splits to isolate, they exhibit shorter path lengths in the decision trees. This approach does not require labeled training data, making it highly effective for identifying unexpected land-use changes \cite{R247}.
  \item \textbf{Geospatial Foundation Models}: The recent introduction of models like IBM-NASA's \textbf{Prithvi-100M} leverages Vision Transformer (ViT) architectures pre-trained on massive Landsat and Sentinel-2 datasets \cite{R204}. These foundation models can be fine-tuned with small training samples to perform complex segmentation and classification tasks, indicating a paradigm shift in spatial analytics \cite{R248}.

\end{itemize}
---

\section{Programmatic API Designs and Decoupled System Architectures}
Modern data engineering separates high-intensity spatial processing from the client presentation layer. This decoupled architecture is essential for web-based GIS applications:
\begin{itemize}
  \item \textbf{OGC Standards}: Web Map Services (WMS) and Web Feature Services (WFS) define standard interfaces for requesting geographic rasters and vector boundaries \cite{R249}.
  \item \textbf{SpatioTemporal Asset Catalog (STAC)}: The STAC specification standardizes geospatial metadata search, allowing pipelines to query diverse data archives ( Landsat, Sentinel, Planet) using a unified query format \cite{R250}.
  \item \textbf{RESTful JSON Exchanges}: By converting heavy spatial data into compressed JSON or GeoJSON vectors, systems can run heavy computations in the cloud (e.g., Earth Engine) and serve the results to lightweight web clients, reducing network bandwidth \cite{R251}.

\end{itemize}
---

\section{Information Engineering Principles (KII Department Alignment)}
This study applies core informatics principles from the Department of Informatics (KII) at the Czech University of Life Sciences Prague (CZU PEF):
\begin{enumerate}
  \item \textbf{ETL Pipeline Orchestration}: Structuring automated Extract, Transform, Load (ETL) workflows to query APIs, clean raw pixels, run machine learning models, and serialize output datasets \cite{R252}.
  \item \textbf{Decision Support Systems (DSS)}: Designing user-centric dashboards that compile complex multi-spectral time-series into clear visual indicators for economic analysts \cite{R253}.
  \item \textbf{Data Quality \& Standardization}: Enforcing data integrity standards, including metadata documentation (ISO 19115) and vector geometry validation (RFC 7946) \cite{R254}.

\end{enumerate}
---

\section{Legal, Ethical, and Security Frameworks}
Satellite remote sensing is subject to international laws, security regulations, and ethical guidelines:
\begin{itemize}
  \item \textbf{Export Control (ITAR \& EAR)}: US export laws regulate spacecraft components and high-resolution imaging technology under ITAR and EAR. However, open-access datasets from civilian platforms (Landsat, Sentinel) are classified as public domain, enabling global academic cooperation \cite{R255}.
  \item \textbf{EU Space Programme (Regulation 2021/696)}: Governs the Copernicus program, establishing a "free, full and open" data policy to foster economic growth and scientific research \cite{R256}.
  \item \textbf{GDPR and Privacy}: High-resolution imagery (under 1 meter) can raise privacy concerns. Using medium-resolution open datasets (10m–30m) in this thesis monitors macro-level activity without tracking individual citizens \cite{R257}.
  \item \textbf{Dual-Use Ethics}: Spatial anomaly detection models can identify industrial or resource changes that have military implications. Researchers must document these dual-use risks and adhere to institutional research ethics \cite{R258}.

\end{itemize}
---

\section{Gap Analysis and Research Positioning}
While existing literature covers remote sensing formulas or economic correlations in isolation, there is a lack of integrated pipelines that ingest, process, detect anomalies, and present results in a unified web-based system. Most academic models remain locked in offline Python environments. This thesis addresses this gap by building a fully automated, reproducible, open-API end-to-end architecture, delivering results through a futuristic web interface.

---

